\chapter{Chlamydia Testing}

\section*{What Is This Test?}
Chlamydia testing detects \textit{Chlamydia trachomatis}, an obligate intracellular gram-negative bacterium and the most commonly reported bacterial STI in the United States, with approximately 1.6 million cases reported annually. \textit{C. trachomatis} has 15 recognized serovars: serovars A--C cause trachoma (the leading infectious cause of blindness worldwide, prevalent in resource-limited settings); serovars D--K cause urogenital infections, inclusion conjunctivitis, and neonatal pneumonia; serovars L1--L3 cause lymphogranuloma venereum (LGV), an invasive disease characterized by genital ulceration, painful inguinal lymphadenopathy (buboes), and proctocolitis (particularly in MSM). Urogenital chlamydia is asymptomatic in 70--85\% of women and 50--75\% of men, facilitating silent transmission. In women, untreated chlamydia is the leading preventable cause of tubal factor infertility, ectopic pregnancy, and chronic pelvic pain due to ascending infection causing salpingitis, PID, and perihepatitis (Fitz-Hugh-Curtis syndrome). In men, complications include epididymitis and urethritis. Neonatal acquisition through an infected birth canal causes inclusion conjunctivitis (5--14 days after birth) and chlamydial pneumonia (4--12 weeks of age, characterized by staccato cough, afebrile course, and eosinophilia).

\section*{Alternative Names}
Chlamydia NAAT; Chlamydia PCR; C. trachomatis Testing; Chlamydia Culture; Chlamydia DNA Probe; CT NAAT; Chlamydia Serology (LGV); Lymphogranuloma Venereum Testing

\section*{Principle}
Nucleic acid amplification tests (NAATs) are the gold standard and the CDC-recommended test for all chlamydia diagnostic testing. These include real-time PCR (Roche cobas CT/NG, Abbott RealTime CT/NG, Hologic Aptima Combo 2, Cepheid GeneXpert CT/NG). Detection of \textit{C. trachomatis} DNA or rRNA using targets in the cryptic plasmid (which is present in approximately 7--10 copies per elementary body) and/or chromosomal targets (\textit{ompA} gene encoding the major outer membrane protein/MOMP). The Hologic Aptima assay targets 23S rRNA, leveraging the high rRNA copy number for improved sensitivity. Sensitivity 95--99\% for urogenital specimens, 90--95\% for rectal and pharyngeal specimens. Specificity >99\%. Vaginal swabs (self-collected or clinician-collected) are the preferred specimen type for screening women; first-catch urine is preferred for screening men. Results available in 1--4 hours. Culture: Specimens inoculated onto cycloheximide-treated McCoy cell monolayers. After 48--72 hours of incubation, \textit{C. trachomatis} inclusions are detected by fluorescein-labeled monoclonal antibody staining. Culture is 100\% specific but has lower sensitivity (60--85\%) compared to NAAT and is labor-intensive, expensive, and technically demanding. It is now reserved for medicolegal cases (sexual abuse), isolates for antimicrobial susceptibility testing, and confirmatory testing when NAAT results are questioned. LGV-specific testing: Genotyping (PCR and sequencing of \textit{ompA} gene) of \textit{C. trachomatis}-positive rectal specimens is required to distinguish LGV serovars (L1--L3) from non-LGV serovars (D--K). LGV genotyping is recommended for all MSM with symptomatic proctitis/proctocolitis and a positive rectal chlamydia NAAT.

\section*{History}
\textit{Chlamydia trachomatis} was first identified by Ludwig Halberstaedter and Stanislaus von Prowazek in 1907 in inclusion bodies from trachoma patients. It was initially misclassified as a virus (TRIC agent: trachoma-inclusion conjunctivitis) until the 1960s when it was recognized as a bacterium based on its possession of both DNA and RNA, a bacterial cell wall, and susceptibility to antibiotics. Cell culture methods developed in the 1960s--1970s aided in understanding the natural history and complications of chlamydia. Antigen detection tests (direct fluorescent antibody/DFA, enzyme immunoassay/EIA) were used in the 1980s--early 1990s before the introduction of NAATs. The first NAAT (PCR, Roche Amplicor) was FDA-cleared in 1993, providing a major advance with improved sensitivity and the ability to use noninvasive specimens (urine). Successive generations of NAATs in the 2000s--2010s achieved higher throughput, faster turnaround, and the ability to test extragenital sites. The CDC's chlamydia screening recommendations (annual screening for sexually active women under 25) were first issued in 1993 and remain core to public health STI control. Reports of LGV proctitis outbreaks among MSM in Europe and North America in the early 2000s prompted renewed attention to LGV and the need for molecular typing.

\section*{Why This Test Is Ordered}
\begin{itemize}
  \item Annual screening of all sexually active women under 25 years of age, and women 25 years and older with risk factors (new or multiple sexual partners, a partner with STI, inconsistent condom use, commercial sex work), for urogenital chlamydia per CDC and USPSTF Grade B recommendation
  \item Screening of all pregnant women at the first prenatal visit; repeat screening in the third trimester for women at ongoing risk. This is critical because chlamydia during pregnancy is associated with preterm birth, premature rupture of membranes, low birth weight, and postpartum endometritis
  \item Screening of MSM at least annually for urogenital, rectal, and pharyngeal chlamydia depending on sexual practices; quarterly (3-monthly) screening for MSM at high risk (multiple partners, methamphetamine use, STI exposure)
  \item Diagnosis of chlamydial urethritis in men with dysuria, urethral discharge, or urethral pruritus
  \item Evaluation of women with mucopurulent cervical discharge, easily induced endocervical bleeding (friability), intermenstrual bleeding, or pelvic/abdominal pain suggestive of cervicitis or PID
  \item Diagnosis of LGV proctitis/proctocolitis in MSM presenting with severe anorectal pain, tenesmus, bloody or purulent rectal discharge, constipation, and inguinal lymphadenopathy (buboes)
  \item Testing of neonates with purulent conjunctivitis (onset 5--14 days after birth) and/or afebrile pneumonia with staccato cough and peripheral eosinophilia (onset 4--12 weeks)
  \item Testing of sexual partners of individuals diagnosed with chlamydia (expedited partner therapy where legal: doxycycline 100 mg PO BID x 7 days prescription dispensed to the index patient for the partner without medical evaluation)
\end{itemize}

\section*{Is This a Primary or Secondary Test}
NAAT is the primary (and only CDC-recommended) diagnostic test for chlamydia at all sites. Culture is a secondary test used primarily in medicolegal settings (child sexual abuse) where the 100\% specificity is required. Serology is NOT recommended for routine diagnosis of acute urogenital chlamydia; it may be used as an adjunct in LGV diagnosis (complement fixation titer >1:64 with compatible clinical syndrome) and, historically, for evaluation of tubal factor infertility.

\section*{How This Test Is Performed}
Vaginal swab (self-collected): The patient is instructed (with visual diagram) to insert the swab 2--3 inches into the vagina, rotate for 10--30 seconds, and withdraw, placing it immediately into the manufacturer's transport tube. Self-collected vaginal swabs are equivalent in sensitivity to clinician-collected endocervical swabs and are the specimen of choice for screening. Endocervical swab: A speculum exam is performed; mucus is removed from the cervical os, and a flocked swab is inserted 1--2 cm into the endocervical canal and rotated against the endocervical epithelium for 15--30 seconds. First-catch urine (FCU): Patient should not have urinated for at least 1 hour; collect the first 10--20 mL of voided urine. Rectal swab: Blindly inserted 2--3 cm into the anal canal, rotated gently against the mucosa. Pharyngeal swab: Swab the tonsillar pillars and posterior oropharynx. All swabs are placed in the manufacturer-specific transport tube. For LGV typing, a separate rectal swab should be sent with a specific request for LGV genotyping (PCR and sequencing of \textit{ompA} gene). First-void urine from men and self-collected vaginal swabs are the highest-yield noninvasive specimens.

\section*{What Preparation Is Needed}
Patients should not urinate for at least 1 hour before first-catch urine collection. Women should not douche, use vaginal creams, or insert suppositories for 24 hours before specimen collection. For optimal yield from endocervical (clinician-collected) swabs, the cervix should not be cleaned with antiseptic (povidone-iodine) before specimen collection. Recent antibiotic therapy reduces NAAT sensitivity; ideally, testing should be deferred for 3 weeks after completion of antimicrobial therapy. For test of cure when indicated, testing should be performed no sooner than 4 weeks after treatment completion (to avoid false positives from residual nonviable nucleic acid).

\section*{Contraindications and Precautions}
NAATs for chlamydia should NOT be used for test of cure earlier than 3--4 weeks after treatment, because residual DNA/rRNA from nonviable organisms can persist and produce false-positive results. Test of cure is NOT routinely recommended after treatment of uncomplicated urogenital chlamydia with a CDC-recommended regimen (doxycycline or azithromycin) because cure rates exceed 95\%. However, test of cure is recommended for: pregnant women (test at 4 weeks after treatment), patients with persistent or recurrent symptoms, patients with suspected nonadherence, and all patients treated for LGV. Chlamydia serology has poor sensitivity and specificity for acute urogenital infection and should NOT be used. Pharyngeal chlamydia is far less common than pharyngeal gonorrhea (<2\% vs. 5--10\% in MSM) and routine pharyngeal screening for chlamydia is not currently recommended (CDC does not recommend routine pharyngeal chlamydia screening). Lymphogranuloma venereum should be considered in any MSM with severe proctitis/proctocolitis, especially with a positive rectal chlamydia NAAT, and specific molecular typing should be requested.

\section*{Risks}
Specimen collection is minimally invasive with negligible risk. Self-collection is completely noninvasive. The primary risks of untreated chlamydia are not from testing but from the disease itself: tubal factor infertility (risk increases with each episode of PID), ectopic pregnancy, and chronic pelvic pain in women; epididymitis and potential contribution to male subfertility in men. Untreated chlamydia during pregnancy is associated with preterm delivery and neonatal chlamydial infection.

\section*{What Happens After the Test}
NAAT results available in 1--4 hours (rapid platforms 90 minutes). Confirmed chlamydia requires treatment with doxycycline 100 mg PO BID for 7 days (preferred first-line, superior to single-dose azithromycin for rectal chlamydia). Azithromycin 1 g PO as a single dose is the alternative (during pregnancy, azithromycin is the recommended agent; doxycycline is contraindicated). Sexual partners from the preceding 60 days should be evaluated, tested, and treated presumptively. EPT is recommended where legal. Patients should abstain from sexual activity for 7 days after single-dose therapy or until completion of a 7-day regimen and resolution of symptoms. All cases must be reported to public health within the required reporting timeframe.

\section*{Understanding Results}

\subsection*{Below Normal}
\begin{itemize}
  \item \textbf{NAAT negative (not detected):} No \textit{C. trachomatis} DNA/rRNA detected at the tested site. High NPV (>99\%)---essentially excludes infection if specimen was adequately collected. However, testing only one anatomic site may miss infection at untested sites (e.g., negative urine NAAT does not exclude rectal chlamydia)
  \item \textbf{Culture negative:} No \textit{C. trachomatis} inclusions detected. High NPV in the appropriate setting but lower sensitivity than NAAT, so false negatives are more common. NAAT should always be the primary test
  \item \textbf{Chlamydia serology negative:} Does not exclude acute chlamydia. Not recommended for routine clinical use
\end{itemize}

\subsection*{Above Normal}
\begin{itemize}
  \item \textbf{NAAT positive for \textit{C. trachomatis}:} Diagnostic of chlamydial infection at the site tested. All patients should be treated. Re-testing at 3 months is recommended for all patients (not a test of cure but to exclude reinfection, which occurs in 10--20\% of treated patients due to resumption of sexual activity with untreated partners)
  \item \textbf{Culture positive:} Definitive diagnosis (100\% specific). Used in medicolegal contexts (sexual abuse) where a false-positive result has grave consequences
  \item \textbf{Rectal NAAT positive for \textit{C. trachomatis} with LGV genotyping positive (serovars L1, L2, L3):} Diagnostic of lymphogranuloma venereum (LGV). Treatment requires doxycycline 100 mg PO BID for 21 days (triple the duration for non-LGV chlamydia). Test of cure is recommended for all LGV cases. Buboes may require aspiration (not incision and drainage) to prevent fistula formation
  \item \textbf{Neonatal conjunctival NAAT positive for \textit{C. trachomatis}:} Inclusion conjunctivitis of the newborn. Oral erythromycin (50 mg/kg/day divided QID x 14 days) or azithromycin (20 mg/kg/day x 3 days) is recommended. Topical antibiotics alone are ineffective; systemic therapy is required to treat nasopharyngeal colonization and prevent subsequent chlamydial pneumonia. The mother and her sexual partner(s) should be tested and treated
\end{itemize}

\section*{Normal Results}
\textbf{\textit{C. trachomatis} not detected} by NAAT at all tested sites. No evidence of chlamydial infection.

\section*{Follow-up Tests}
\begin{itemize}
  \item \textbf{Re-testing for chlamydia (NAAT):} All patients diagnosed with chlamydia should be re-tested approximately 3 months after treatment, regardless of whether symptoms are present or partners were treated. This is NOT a test of cure but screening for reinfection, which occurs in 10--20\% within 3--6 months
  \item \textbf{Test of cure (NAAT) at 4 weeks post-treatment:} Indicated for pregnant women, LGV, persistent symptoms, or suspected nonadherence. Must be done at least 4 weeks after treatment completion to avoid false positives
  \item \textit{\textbf{Neisseria gonorrhoeae}} \textbf{NAAT from the same specimens (dual CT/NG NAAT):} Co-testing for gonorrhea at all sites. Co-infection rates in chlamydia-positive patients are 5--15\%
  \item \textbf{HIV testing, syphilis serology:} Recommended for all patients diagnosed with chlamydia given shared risk factors
  \item \textbf{LGV genotyping (PCR and sequencing of \textit{ompA}):} For all rectal chlamydia-positive specimens from MSM with signs/symptoms of proctitis or from those in regions where LGV is known to circulate
\end{itemize}

\section*{References}
\begin{enumerate}
  \item Workowski KA, Bachmann LH, Chan PA, et al. Sexually transmitted infections treatment guidelines, 2021. MMWR Recomm Rep. 2021;70(4):1-187.
  \item CDC. Recommendations for the laboratory-based detection of \textit{Chlamydia trachomatis} and \textit{Neisseria gonorrhoeae} --- 2014. MMWR Recomm Rep. 2014;63(RR-02):1-19.
  \item Geisler WM. Duration of untreated, uncomplicated \textit{Chlamydia trachomatis} genital infection and factors associated with clearance. J Infect Dis. 2011;204(Suppl 2):S134-S155.
  \item Haggerty CL, Gottlieb SL, Taylor BD, et al. Risk of sequelae after \textit{Chlamydia trachomatis} genital infection in women. J Infect Dis. 2010;201(Suppl 2):S134-S155.
  \item Stoner BP, Cohen SE. Lymphogranuloma venereum 2015: clinical presentation, diagnosis, and treatment. Clin Infect Dis. 2015;61(Suppl 8):S865-S873.
\end{enumerate}

\clearpage
